\documentclass[11pt,a4paper]{article}

\usepackage[margin=1in]{geometry}
\usepackage[T1]{fontenc}
\usepackage[utf8]{inputenc}
\usepackage{lmodern}
\usepackage[none]{hyphenat}
\usepackage{graphicx}
\usepackage{booktabs}
\usepackage{tabularx}
\usepackage{longtable}
\usepackage{array}
\usepackage{hyperref}
\usepackage{xcolor}
\usepackage{caption}
\usepackage{float}
\usepackage{pdflscape}

\hypersetup{
  colorlinks=true,
  linkcolor=blue!60!black,
  urlcolor=blue!60!black,
  citecolor=blue!60!black,
  pdftitle={Final Project Report - STEM Learning Platform with GenAI},
  pdfauthor={Group 8}
}

\graphicspath{
  {../docs/mermaid-png/root/}
  {../docs/mermaid-png/web/}
  {../docs/mermaid-png/backend/}
  {../docs/mermaid-png/supabase/}
}

\setlength{\parindent}{0pt}
\setlength{\parskip}{0.65em}
\renewcommand{\arraystretch}{1.2}
\hyphenpenalty=10000
\exhyphenpenalty=10000

\newcommand{\ProjectTitle}{AI STEM Learning Platform}
\newcommand{\ReportTitle}{Final Project Report}
\newcommand{\CourseName}{COMP3122 Information Systems Development}
\newcommand{\GroupNumber}{8}
\newcommand{\PreparedBy}{Group \#\GroupNumber}
\newcommand{\ReportDate}{April 2026}
\newcommand{\LiveProductUrl}{\href{https://ai-stem-learning-platform-group-8.vercel.app/}{\nolinkurl{ai-stem-learning-platform-group-8.vercel.app}}}
\newcommand{\DemoVideoUrl}{\href{https://youtu.be/xsq_rRhuclY}{\nolinkurl{youtu.be/xsq_rRhuclY}}}
\newcommand{\PrimaryRepoUrl}{\href{https://github.com/HaochenFa/AI-Learning-Platform-for-STEM}{\nolinkurl{HaochenFa/AI-Learning-Platform-for-STEM}}}
\newcommand{\SubmissionRepoUrl}{\href{https://github.com/comp3122-2526sem2/groupproject-team_8}{\nolinkurl{comp3122-2526sem2/groupproject-team_8}}}
\newcommand{\PageFigureWidth}{\dimexpr\paperwidth-1.25in\relax}
\newcommand{\FullPageWidthGraphic}[1]{%
  \makebox[\textwidth][c]{%
    \includegraphics[
      width=\PageFigureWidth,
      height=0.74\textheight,
      keepaspectratio
    ]{#1}%
  }%
}
\begin{document}
\hypersetup{pageanchor=false}
\pagenumbering{gobble}

\begin{titlepage}
\centering
\vspace*{2cm}
{\Large \CourseName \par}
\vspace{0.4cm}
{\large Group \#\GroupNumber \par}
\vspace{1cm}
{\Huge \bfseries \ReportTitle \par}
\vspace{0.8cm}
{\LARGE \ProjectTitle \par}
\vspace{1.5cm}
\begin{tabularx}{1\textwidth}{>{\bfseries\raggedright\arraybackslash}p{4.4cm}@{\hspace{1cm}}X}
Prepared by & \PreparedBy \\
Date & \ReportDate \\
Repo snapshot & 2026-04-05 \\
Live product & \LiveProductUrl \\
Demo video & \DemoVideoUrl \\
Primary repository & \PrimaryRepoUrl \\
Submission repository & \SubmissionRepoUrl \\
\end{tabularx}
\vspace{1cm}

{\large \bfseries Group Members\par}
\vspace{0.35cm}
\begin{tabularx}{0.9\textwidth}{>{\bfseries\raggedright\arraybackslash}p{4.2cm}@{\hspace{1.2cm}}>{\raggedright\arraybackslash}X}
\toprule
Name & Student ID \\
\midrule
\mbox{FA, Haochen} & \mbox{24113347D} \\
\mbox{WONG Ka Hei} & \mbox{23118174D} \\
\mbox{KWAN Kai Man} & \mbox{23118556D} \\
\mbox{LI, Zhengyang} & \mbox{24099899D} \\
\bottomrule
\end{tabularx}
\vfill
{\large
This report documents the current implemented state of the repository, including the
web application, Python backend, Supabase schema and Edge Functions, testing strategy,
deployment topology, and role-based user workflows.
\par}
\end{titlepage}

\clearpage
\hypersetup{pageanchor=true}
\pagenumbering{arabic}
\tableofcontents
\newpage
\raggedright

\section{Executive Summary}

This project delivers a production-oriented STEM learning platform that places Generative AI inside
the real classroom operating model rather than exposing it as a detached chatbot. Its core idea is
simple but powerful: teachers should be able to transform their own teaching materials into a
structured \emph{Course Blueprint}, and that blueprint should govern what students study, what the
AI can generate, and what the teacher can review afterwards.

That blueprint-centered loop gives the product its practical value. Teachers upload materials, the
system processes and structures them, AI helps generate and refine the class blueprint, students
learn through blueprint-grounded chat and activities, and the resulting student work flows back into
review, class intelligence, and adaptive teaching support. In other words, AI is used here to
support the full teaching cycle from preparation to intervention, not just to answer questions on
demand.

The implemented platform serves three concrete audiences. Teachers receive management and decision
surfaces for class creation, material curation, blueprint publishing, assignment authoring,
submission review, chat monitoring, and teaching insight interpretation. Students receive a separate
guided workspace for join-by-code class access, published blueprint study, always-on class chat, and
tracked assignment completion. Guests receive a sandboxed but realistic evaluation path with
anonymous entry, role switching, reset, and data isolation, making the product demonstrable without
touching real classroom data.

Operationally, the system is implemented across three cooperating runtime surfaces: a \textbf{Next.js
16} web application for role-aware user experience and server actions, a \textbf{Python FastAPI}
backend for AI orchestration and workflow-heavy service logic, and a \textbf{Supabase} project for
authentication, PostgreSQL persistence, row-level security, storage, and background-job support.
This split gives the platform clearer operational boundaries and allows AI orchestration,
persistence, and lifecycle concerns to be treated as first-class system responsibilities.

The current repository also provides strong evidence of implementation maturity:
\begin{itemize}
  \item \textbf{Live product:} the deployed system is available at \LiveProductUrl.
  \item \textbf{Repository scale:} the current snapshot contains 411 commits, including 92 merge
  commits, alongside 23 active database migrations.
  \item \textbf{Quality signals:} the codebase currently includes 60 web test files, 17 backend
  test files, 9 Playwright end-to-end specifications, and a documented deployment topology spanning
  the web app, Python backend, and Supabase services.
\end{itemize}

Taken together, the current snapshot represents a coordinated educational product with credible
implementation depth and clear operational boundaries, not a toy prototype or a one-off AI demo.

\section{Project Background and Motivation}

\subsection{Course and problem context}

The course brief frames the project in the broader language of \emph{Education 4.0}, personalized
learning, and AI-supported teaching. That framing is useful, but the repository shows that the real
product challenge is more concrete: classroom AI must do more than answer questions. Teachers need a
system that can absorb their materials, structure them into trustworthy course knowledge, turn that
knowledge into student-facing activities, and still leave the teacher in editorial control.

In practice, the pain points are operational as much as pedagogical. A teacher does not only need
``an AI''; they need help with lesson preparation, activity creation, assignment delivery, review,
and follow-up instruction. Students, in turn, do not benefit much from a generic chatbot that has no
understanding of what their class is actually studying. They need a guided workspace grounded in the
same materials, objectives, and assignments that the teacher is using.

This project responds to that problem by treating the class itself as the unit of AI orchestration.
Uploaded teaching materials become the basis for a structured Course Blueprint, the blueprint becomes
the source of truth for chat and activities, and the resulting student work flows back into review
and teacher intelligence surfaces. That is the central product idea that differentiates the system
from a thin AI wrapper or a one-off classroom demo.

\subsection{Project goals}

Based on the implemented system and supporting documentation, the project goals can be expressed as
four connected product objectives rather than a flat checklist.

\paragraph*{Structured knowledge creation.}
The first goal is to transform raw classroom materials into reusable, class-specific knowledge. The
implemented platform does this through a materials library, asynchronous processing, and AI-assisted
blueprint generation so that uploaded files become an organised teaching asset instead of a passive
attachment store.

\paragraph*{Teacher control and trust.}
The second goal is to keep instructors in control of what AI produces. The blueprint workflow is
therefore not instant-publish generation: teachers edit the draft, approve it into an overview
state, and publish it only when it is ready for students. The same idea extends to quiz and
flashcard drafts, which are editable before assignment.

\paragraph*{Guided student learning.}
The third goal is to give students a classroom workspace that feels specific to their class rather
than generic to the internet. The implemented student experience centres on the published blueprint,
always-on class chat, assignable quiz/flashcard/chat activities, and visible progress states such as
assigned, in progress, submitted, and reviewed.

\paragraph*{Operational credibility.}
The fourth goal is to demonstrate that the platform is credible as a real software product, not only
as a concept. The current repository therefore includes a dedicated backend orchestration layer,
queue-backed material processing, deployment documentation, automated tests across web and backend
surfaces, and a sandboxed guest evaluation mode that allows realistic product walkthroughs without
touching live classroom data.

\subsection{What makes this project distinctive}

What makes the project distinctive is not any single AI endpoint, but the way several design choices
combine into a more complete classroom system.

\paragraph*{Blueprint as the pedagogical control layer.}
The Course Blueprint is not decorative metadata. It is the product's central teaching object: it is
generated from class materials, edited by the teacher, versioned through draft/overview/published
states, and then reused to ground downstream student experiences. This adds a meaningful editorial
checkpoint that most GenAI classroom demos do not have.

\paragraph*{A full classroom loop rather than isolated generation.}
The system supports an end-to-end workflow from teacher setup to student activity to teacher review.
That means AI is used not just to generate content, but to support assignment creation, submission
handling, feedback, chat monitoring, class insights, and teaching briefs. The product therefore
models classroom operations, not only content synthesis.

\paragraph*{A deliberate AI orchestration boundary.}
All AI generation routes through a dedicated Python FastAPI service rather than ad hoc provider calls
from the web layer. This gives the system a clearer separation of concerns, consistent response
contracts, and a stronger basis for provider routing, validation, and future extension.

\paragraph*{Production-minded background and persistence design.}
Material processing is asynchronous and queue-backed, analytics and teaching briefs are snapshot
oriented, and blueprint publication is treated as a state transition rather than a transient prompt
result. These choices make the product feel much closer to a deployable platform than to a hackathon
prototype.

\paragraph*{An authentic guest evaluation path.}
Guest mode is implemented as a real sandbox classroom with anonymous authentication, seeded demo
data, role switching, reset, expiry, and cleanup. That is a significantly stronger demonstration
surface than a static landing page tour or a pre-recorded slideshow of product screens.

\section{System Overview and Implemented Features}

\subsection{High-level system picture}

Figure~\ref{fig:system-overview} shows the top-level structure of the current platform.

\begin{figure}[H]
  \centering
  \FullPageWidthGraphic{README-diagram-01.png}
  \caption{Current high-level platform architecture.}
  \label{fig:system-overview}
\end{figure}

At a high level, teacher, student, and guest users all enter through the web application.
The web layer handles UI rendering and user flow orchestration, while the backend is responsible
for AI-heavy logic and the Supabase layer provides persistence, storage, security, and async jobs.

\subsection{Primary user groups}

\begin{table}[H]
\centering
\caption{Primary user groups in the current release}
\begin{tabularx}{\textwidth}{>{\bfseries}p{3cm}X}
\toprule
Role & Implemented responsibilities \\
\midrule
Instructor & Create and manage classes, upload materials, curate and publish blueprints, generate and assign activities, review submissions, monitor chat behaviour, inspect insights, use the adaptive teaching brief, and preview the student journey from within the class workspace. \\
Student & Join classes with a join code, study the published blueprint, use always-on class chat, complete assigned chat/quiz/flashcard activities, and track current, upcoming, completed, and reviewed work from the student dashboard. \\
Guest & Enter a sandbox classroom through anonymous auth, explore seeded teacher and student workflows through role switching, reset the demo when needed, and evaluate the product without affecting real user data. \\
\bottomrule
\end{tabularx}
\end{table}

These user groups are not theoretical labels; they map directly to different product surfaces and
navigation models in the implementation.

\paragraph*{Instructor.}
The instructor experience is creation-first and decision-heavy. Teachers land in a dedicated
dashboard, see their owned and assistant classes, and move into a class shell where they can upload
materials, generate blueprints, create assignments, monitor student chats, and interpret class
insights. The product also supports preview-as-student mode, which lets a teacher inspect the
student-facing experience without leaving the classroom context.

\paragraph*{Student.}
The student experience is intentionally simpler and more guided. Students enter through a separate
dashboard, join classes through join codes, and then work inside a class workspace organised around
the published blueprint, always-on chat, and assigned activities. Their dashboard and class pages
emphasise due work, progress state, and reviewed feedback rather than creation or analytics tools.

\paragraph*{Guest.}
The guest experience is designed for realistic evaluation rather than passive observation. Guests
authenticate anonymously, receive a seeded sandbox classroom, and can switch between teacher and
student perspectives, reset the environment, or convert to a real signup path. Middleware and
sandbox state keep the guest confined to the sandbox classroom so the demo remains safe and
operationally believable.

In addition to these primary personas, the implementation also supports a class-scoped
\textbf{teaching assistant (TA)} authorization role. TA is not exposed as a separate permanent
signup account type; instead, it exists as an enrollment-level role that can access the same
teacher-oriented class workflows where route and service logic explicitly permit \texttt{teacher}
or \texttt{ta}.

\subsection{Feature set}

Rather than presenting the system as a single AI chatbot, the current release delivers a set of
user-facing classroom capabilities that work together as a guided teaching and learning platform.
The most important design decision is that AI is always embedded inside a teacher-controlled class
workflow: content is grounded in uploaded materials, reviewed through a blueprint, released through
activities, and then interpreted through feedback and insights. The main implemented feature groups
are summarised below.

\subsubsection*{Role-Based Access and Entry}

The platform supports distinct entry paths for teachers and students rather than a single generic
login followed by ad hoc navigation. Permanent users sign in with email and password, complete
email verification, and then land in role-specific dashboards that match their responsibilities.
Teachers enter management-oriented flows for classes, materials, blueprints, assignments, and
insights; students enter a simpler workspace centred on class participation and assigned work.

This access model is implemented as a shared authentication surface rather than duplicated forms.
The same auth component handles sign-in, sign-up, and password recovery in both the landing-page
modal and dedicated auth pages, while server-side confirmation and recovery routing supports secure
email verification and password reset completion. The result is a cleaner product experience for
users and a more maintainable authentication boundary for the system.

\subsubsection*{Classroom Setup and Participation}

For teachers, the first meaningful product workflow is class setup. A teacher can create a class,
define its basic metadata, and immediately receive a join code that can be shared with students.
Students can then join an existing class through a join-by-code flow and are placed directly into a
class-scoped workspace with the right navigation, permissions, and learning surfaces for their
role.

This capability is more than a simple CRUD screen. Class creation automatically enrolls the
creating teacher, join codes are validated through a dedicated backend workflow, and join operations
are designed to be safely repeatable rather than fragile one-time mutations. Once inside a class,
the platform keeps users in a consistent class shell, and teachers can even preview the student
experience in read-only mode without leaving the same classroom context.

\subsubsection*{Material Ingestion and Blueprint Publishing}

The core differentiator of the platform is its material-to-blueprint workflow. Teachers can upload
PDF, DOCX, and PPTX teaching materials into a class library, monitor their processing status, and
manage them through preview, download, and delete actions. This turns the class workspace into a
real content hub rather than a thin wrapper around a text prompt box.

Once materials are available, the platform uses them to generate a structured Course Blueprint that
becomes the class source of truth. The generated blueprint is not exposed to students immediately:
teachers edit the draft, adjust topics, objectives, and prerequisite structure, approve it into an
overview state, and then publish it when ready. This creates a clear editorial checkpoint and makes
the AI output reviewable rather than automatic.

The implementation depth behind this feature is substantial. Material ingestion is asynchronous,
with extraction, chunking, and embedding handled in the background so the system does not pretend
that an uploaded file is instantly AI-ready. Blueprint generation produces structured payloads
rather than free-form prose, while publication is owner-controlled and version-aware, allowing a
teacher to create a new draft from the currently published blueprint without losing the stable
student-facing version until the revised content is explicitly promoted.

\subsubsection*{AI Learning Activities for Students}

After a blueprint is published, the platform turns that curated course knowledge into multiple
learning experiences for students. Students can use an always-on class chat for open practice, and
teachers can create formal chat, quiz, and flashcard assignments from the class blueprint. This
gives the product a broader learning model than a one-dimensional question-answer interface.

Quiz and flashcard activities are generated as editable drafts, allowing teachers to refine the AI
output before publishing it and assigning it to the whole class. Chat assignments add
teacher-authored instructions to a guided conversation flow, while student dashboards and class
views track assigned, in-progress, submitted, and reviewed states. Quiz flows support repeated
attempts and score tracking, flashcard sessions capture known-versus-needs-review outcomes, and
chat assignments preserve both the conversation transcript and student reflection for later review.

The student experience is also visually richer than plain text output. In both chat and data-rich
contexts, the system can attach generative canvas responses such as charts and diagrams when a
concept is better explained visually. This is especially valuable in STEM learning, where ideas
such as processes, relationships, and trends are often easier to understand when rendered instead
of only described.

\subsubsection*{Teacher Review and Instructional Intelligence}

The platform does not stop at content generation or assignment delivery. Teachers can open dedicated
review screens for chat, quiz, and flashcard assignments to inspect what students actually
submitted, check completion status across the class, and record feedback. This shifts the product
from ``AI makes content'' to ``AI supports a full teaching cycle''.

The review experience is implemented with meaningful depth. Teachers can inspect quiz attempts and
answers, read chat transcripts, review flashcard session outcomes, override or adjust scores where
appropriate, and save comments plus highlights back to the student record. In parallel, a
read-only chat monitor lets teachers inspect students' always-on class chats for coaching and
support without interfering with the student's active workspace.

On top of review, the platform provides a teacher intelligence layer. The class intelligence
dashboard aggregates quiz performance, topic outcomes, Bloom-level coverage, engagement patterns,
student drill-downs, and intervention suggestions. A natural-language data-query panel lets a
teacher ask questions about class performance and receive generated charts, while the adaptive
teaching brief condenses the current class state into a memo-style summary with attention items,
students to watch, likely misconceptions, strongest recommended action, and suggested next steps.
These features are backed by persisted snapshots so teacher-facing summaries remain stable across
reloads and can be refreshed intentionally instead of regenerating unpredictably on every visit.

\subsubsection*{Guest Sandbox Evaluation Mode}

Guest mode is one of the most distinctive parts of the release because it is implemented as a real
interactive sandbox rather than a fake product tour. A visitor can enter from the landing page,
authenticate anonymously, and receive a cloned classroom seeded with demo materials, blueprint
data, activities, analytics, and chat history. This allows stakeholders to experience realistic
teacher and student flows without touching live user data.

The sandbox experience includes role switching between teacher and student perspectives, reset
capability for restarting the demo, and clean handoff into permanent signup when a guest wants to
become a real user. It is also operationally disciplined: guest sessions are isolated through
sandbox-scoped data, confined to allowed routes, limited by global session quotas, and expired
after inactivity or maximum lifetime. A dedicated cleanup workflow then removes guest data,
associated storage objects, and anonymous auth users so the feature remains sustainable in a
hosted environment rather than becoming a one-off demo hack.

\subsection{Blueprint-centered value flow}

The most important end-to-end educational flow in the current release is shown in
Figure~\ref{fig:value-flow}. Unlike the lighter overview used in the repository README, this
workflow captures the implemented control points that make the platform educationally credible:
background material readiness, teacher-governed blueprint publication, blueprint-gated activity
delivery, persisted student work, and teacher-facing instructional response.

\begin{landscape}
\begin{figure}[p]
  \centering
  \includegraphics[
    width=0.96\linewidth,
    height=0.82\textheight,
    keepaspectratio
  ]{ARCHITECTURE-diagram-13.png}
  \caption{Implemented blueprint-governed classroom workflow from material ingestion to student activity, review, and teaching intervention.}
  \label{fig:value-flow}
\end{figure}
\end{landscape}

\paragraph*{Material ingestion and readiness.}
The flow begins with teacher-owned classroom content rather than with a prompt box. When an
instructor uploads PDF, DOCX, or PPTX materials, the web application stores the file and material
record immediately, but it does not pretend the content is AI-ready at once. Instead, the Next.js
server action dispatches a background processing request to the Python backend, which enqueues the
work and can trigger the Supabase \texttt{material-worker} Edge Function. That worker extracts text,
chunks it, generates embeddings, and persists retrieval-ready material context. This is an important
implementation choice because it turns uploaded files into operational class knowledge through a
real asynchronous pipeline, with clear \texttt{processing}, \texttt{ready}, and \texttt{failed}
states, rather than through a fragile synchronous shortcut.

\paragraph*{Blueprint generation and publication control.}
Once materials are ready, the platform retrieves indexed class context and sends it to the Python
backend for structured blueprint generation. The backend does not return informal prose; it returns a
validated payload containing summary, topics, objectives, prerequisite structure, and related
teaching metadata. The web layer then persists that output both as a canonical blueprint snapshot and
as normalized relational records for topics and objectives. Crucially, this generated blueprint is
not student-facing by default. Teachers edit the draft, validate topic relationships, approve it into
an overview state, and publish it only when they are satisfied with the course structure. The owner
controls approval and publication, and the implementation also supports creating a new draft from the
published version so the live student-facing blueprint remains stable until a revised version is
explicitly promoted.

\paragraph*{Blueprint-governed learning delivery.}
After publication, the blueprint becomes the governing instructional artifact for downstream AI
features. Open-practice class chat, quiz generation, flashcard generation, and teacher-created chat
assignments all depend on the published blueprint rather than bypassing it. This is where the
product's core idea becomes concrete: students do not interact with a generic model that improvises
outside the class context; they interact with AI features constrained by approved course knowledge and
retrieved material evidence. The implementation distinguishes between activity types in a sensible
way. Quiz and flashcard activities are generated as editable drafts so teachers can refine the AI
output before release, while chat assignments publish directly because the teacher is already
supplying the assignment framing and instructions in one step.

\paragraph*{Student work and teacher review.}
The workflow then moves from content delivery into accountable classroom participation. When a
teacher assigns an activity, the system creates assignment records for enrolled students and tracks
their status across states such as assigned, in progress, submitted, and reviewed. Students study
the published blueprint, use always-on grounded class chat, and complete quiz, flashcard, and chat
assignments inside a class-scoped workspace. Their attempts, transcripts, reflections, and feedback
artifacts are persisted rather than discarded after generation. This persistence matters because it
allows teachers to inspect actual student work through review screens, save comments and highlights,
check quiz attempts and outcomes, and monitor always-on class chat activity without interrupting the
student's learning flow.

\paragraph*{Insights and adaptive teaching response.}
The loop closes when student activity becomes instructional intelligence rather than raw interaction
data. Persisted submissions and chat outcomes feed analytics snapshots, visual data-query responses,
risk indicators, topic-level performance summaries, and intervention suggestions. On top of that, the
adaptive teaching brief condenses current class evidence into a memo-style teaching response with
attention items, misconceptions, students to watch, strongest recommended action, and suggested next
steps. This means the platform is not merely generating classroom content; it is helping the teacher
interpret what happened and decide what to do next.

Taken together, this flow is the clearest expression of the product's value. The platform does not
stop at AI-assisted content creation, and it does not treat classroom interaction as disposable chat
traffic. Instead, it implements a governed instructional cycle in which teacher materials become
structured class knowledge, that knowledge governs student activity, student activity produces
reviewable evidence, and that evidence informs the teacher's next instructional move.

\section{System Design}

\subsection{Technology stack}

\begin{table}[H]
\centering
\caption{Core implementation stack}
\begin{tabularx}{\textwidth}{>{\bfseries}p{3.5cm}X}
\toprule
Layer & Technology \\
\midrule
Frontend & Next.js 16 App Router, React 19, TypeScript, Tailwind CSS 4, shared UI primitives, Motion for React, Lucide icons \\
Backend & Python FastAPI, httpx, structured Pydantic schemas, provider adapters \\
Database and auth & Supabase PostgreSQL, Supabase Auth, Row Level Security \\
Storage and async processing & Supabase Storage, pgmq queue, pg\_cron, Supabase Edge Functions \\
Testing & Vitest, React Testing Library, Python unittest, Playwright \\
Deployment model & Separate web and backend deployments plus hosted Supabase \\
\bottomrule
\end{tabularx}
\end{table}

\subsection{Architectural decomposition}

The platform is deliberately partitioned into four top-level areas that map cleanly to runtime
responsibility:
\begin{itemize}
  \item \texttt{web/}: user-facing routes, UI components, middleware, and server actions;
  \item \texttt{backend/}: AI orchestration, chat workspace logic, class workflows, analytics;
  \item \texttt{supabase/}: SQL migrations, local config, and Edge Functions;
  \item \texttt{tests/}: Playwright end-to-end coverage running against a deployed target URL.
\end{itemize}

This separation is important because each layer owns the work it is best positioned to execute. The
web application owns rendering, navigation, and server-action orchestration. The Python backend
owns AI routing, chat-workspace behavior, analytics, and workflow-heavy service logic. Supabase
owns authentication, PostgreSQL persistence, row-level security, storage, and queue-backed
background operations. The result is a clearer product boundary between interface concerns,
orchestration concerns, and stateful platform concerns.

\subsection{Web application design}

The web application delivers the platform's role-aware operating surface. It separates teacher and
student experience from the first navigation step, keeps users inside a persistent class shell once
they enter a classroom, and uses server actions plus middleware to make state transitions feel
consistent across the product.

Important web-layer characteristics include:
\begin{itemize}
  \item role-specific dashboards at \texttt{/teacher/dashboard} and \texttt{/student/dashboard};
  \item a home-first auth model via a shared \texttt{AuthSurface} used in the landing-page modal
  and dedicated routes such as \texttt{/login}, \texttt{/register}, and \texttt{/forgot-password};
  \item middleware-enforced redirects for unauthenticated, unverified, and expired guest sessions;
  \item server actions for class creation, join-by-code, materials, assignments, and settings;
  \item a shared design system with centralized icons, tokens, and motion presets.
\end{itemize}

On the teacher side, the class overview acts as the operational control hub for materials,
assignments, teaching brief, and class monitoring. On the student side, the class workspace is
organized around focused learning surfaces for chat, assignments, and the published blueprint
rather than teacher-oriented management tools. A teacher can also preview the student experience
from within the same classroom context, which is a strong product-control detail because it lets an
instructor inspect the student-facing flow without changing accounts or leaving the class shell.

The authentication surface is also more disciplined than a standalone login form. The platform uses
the same shared auth presenter for landing-page modal entry and fallback dedicated routes, supports
server-side email confirmation and password-recovery callbacks through \texttt{/auth/confirm}, and
completes recovery through a dedicated \texttt{/reset-password} flow. Reusing the same auth
component and redirect model across public entry points and email-link return paths reduces UI drift
and keeps authentication behavior easier to maintain.

\subsection{Python backend design}

The Python backend serves as the platform's AI and workflow control plane. That is a significant
architectural choice because it keeps prompt construction, provider routing, response normalization,
and service-to-service validation out of the web layer and behind a dedicated FastAPI boundary.

The backend provides endpoints and domain logic for:
\begin{itemize}
  \item blueprint generation;
  \item quiz generation;
  \item flashcard generation;
  \item grounded chat;
  \item chat workspace sessions and message persistence;
  \item visual canvas generation;
  \item class create and class join;
  \item materials dispatch;
  \item analytics, teaching brief, and natural-language data queries.
\end{itemize}

All normal backend responses follow a consistent envelope:
\begin{center}
\texttt{\{ ok, data, error, meta \}}
\end{center}

That contract gives the web adapters a uniform integration surface for success, failure, and
request metadata, while allowing the backend to evolve internal provider logic, prompt structure,
and workflow composition without forcing parallel changes across many frontend call sites.

\subsection{AI orchestration design}

AI orchestration is environment-driven and intentionally centralized in the backend. The platform
supports three pluggable providers --- OpenRouter, OpenAI, and Gemini --- and the provider layer
resolves an ordered attempt list from the configured default plus the providers that are actually
usable in the current environment.

This design has several advantages:
\begin{itemize}
  \item the web app remains provider-agnostic;
  \item deployment environments can switch provider preference through configuration;
  \item provider failures can be isolated inside the backend;
  \item request deadlines are enforced centrally rather than inconsistently per route.
\end{itemize}

The orchestration layer is also more robust than a simple ``try another API key'' fallback. It
applies a single wall-clock deadline across provider attempts, recomputes the remaining budget
before each call, and uses the same ordered fallback approach for both generation and embeddings.
Operationally, that means a deployment can prefer one provider by configuration, degrade to another
when necessary, and still preserve a stable frontend contract.

The AI surface is broader than content generation alone. In addition to blueprint, quiz, and
flashcard generation, the backend also supports grounded chat, analytics narratives, teaching brief
synthesis, embeddings, and visual canvas specifications. This gives the platform a unified AI layer
for both student learning workflows and teacher decision-support workflows.

\subsection{Chat workspace and memory design}

The chat workspace is one of the platform's most technically mature subsystems because it is designed
as a long-lived classroom conversation system rather than a disposable prompt-response loop. The
backend supports participant listing for teacher monitoring, session creation, rename and archive
operations, paginated message history, and persisted assistant responses so a class conversation can
continue over time instead of resetting on every visit.

\begin{figure}[H]
  \centering
  \FullPageWidthGraphic{ARCHITECTURE-diagram-08.png}
  \caption{Chat workspace flow with grounded context retrieval and persisted assistant responses.}
  \label{fig:chat-workspace}
\end{figure}

The grounding path is equally important. Workspace chat is constrained by published blueprint context
and retrieved material chunks rather than a generic free-form prompt, which keeps class chat tied to
approved course knowledge. For longer conversations, the system uses summary-based compaction and
context-window management instead of replaying the full transcript indefinitely. That combination of
persistence, grounding, and compaction makes the chat surface much closer to a real classroom tool
than to a one-shot chatbot widget.

\subsection{Material processing pipeline}

Material ingestion is implemented as a real asynchronous processing pipeline. Uploading a file creates
the material record and stores the source asset immediately, but AI-readiness is produced through a
separate dispatch-and-worker path rather than by blocking the teacher's request until extraction and
embedding complete.

\begin{figure}[H]
  \centering
  \FullPageWidthGraphic{DEPLOYMENT-diagram-02.png}
  \caption{Asynchronous material-processing pipeline.}
  \label{fig:material-pipeline}
\end{figure}

The implemented flow is more substantial than a background task checkbox. The web layer persists the
material, the Python backend calls an \texttt{enqueue\_material\_job} RPC, PostgreSQL writes the job
into a \texttt{pgmq}-backed queue, and the Supabase \texttt{material-worker} Edge Function extracts
text, chunks content, generates embeddings, and writes retrieval-ready records plus status updates.
The worker can be awakened immediately after dispatch or picked up through cron-backed execution,
which gives the system both responsiveness and resilience.

This design matters because material processing is exactly the kind of workload that becomes fragile
when treated as a synchronous request. Queue-backed execution improves retry behavior, makes
\texttt{processing}, \texttt{ready}, and \texttt{failed} states explicit in the product, and turns raw
classroom files into indexed instructional context rather than opaque attachments.

\subsection{Analytics and adaptive teaching support}

Teacher-facing intelligence is a meaningful part of the platform design, not a thin reporting add-on.
The analytics layer turns persisted submissions, activity outcomes, and chat-derived signals into a
decision-support surface for instructors.

The class intelligence dashboard includes:
\begin{itemize}
  \item summary cards for student count, average score, completion rate, and risk indicators;
  \item topic-level performance charts;
  \item Bloom-level breakdown visualization;
  \item student engagement views and student overview tables;
  \item AI-generated narrative interpretation and intervention suggestions;
  \item natural-language data query support for generating chart specifications.
\end{itemize}

These features are backed by persisted insight snapshots and caching logic so the teacher is not
looking at a completely re-generated story on every reload. The adaptive teaching brief complements
the larger dashboard by delivering a memo-style summary directly on the class overview page, with
recommended next steps, likely misconceptions, students to watch, and follow-up activity guidance.
Together, the dashboard and brief show that the product is designed to help teachers interpret class
evidence, not only generate student content.

\subsection{Security and access-control design}

Security is layered across all three runtime surfaces rather than delegated to any single framework:
\begin{itemize}
  \item \textbf{web layer}: route protection, email verification enforcement, and guest route confinement;
  \item \textbf{backend layer}: service authentication, user bearer-token verification, and guest AI guardrails;
  \item \textbf{Supabase layer}: row-level security, teacher/enrollment policies, sandbox-aware isolation, and cleanup-capable SQL services.
\end{itemize}

Several implementation details show that this hardening model is intentional rather than cosmetic:
\begin{itemize}
  \item permanent users are limited to teacher or student account types;
  \item \texttt{profiles.account\_type} is immutable after signup;
  \item class-scoped teacher permissions can also be granted through the enrollment-level
  \texttt{ta} role without turning TA into a separate permanent account type;
  \item guest sessions are scoped by \texttt{sandbox\_id};
  \item cursor inputs to PostgREST are validated before interpolation;
  \item service-to-service auth between web and backend uses a shared backend API key.
\end{itemize}

The combined effect is important. Route guards shape who can enter a workflow, backend verification
checks who is calling sensitive service endpoints, and RLS plus sandbox-aware policies determine what
data each actor can actually read or mutate. This is the right design for a platform that mixes real
classroom data, AI workflows, and anonymous guest sandboxes in the same hosted system.

\subsection{Guest sandbox design}

Guest mode is one of the product's strongest differentiators because it is implemented as a real
interactive sandbox rather than as a scripted tour. The platform creates an anonymous Supabase auth
session, provisions a \texttt{guest\_sandboxes} record, and clones seeded classroom content into normal
application tables under a sandbox-specific \texttt{sandbox\_id}. Guests can then switch between teacher
and student views, reset the demo, and experience realistic product behavior without touching real
user data.

The quota and lifecycle model combines Supabase-side and backend-side guardrails:
\begin{itemize}
  \item a Supabase-managed \texttt{guest\_session\_quota} control table with RPCs that gate
  sandbox creation and quota release;
  \item default admission caps of 60 active guest sandboxes globally and 20 newly created
  sandboxes per hour;
  \item backend per-feature guest AI limits plus an \texttt{asyncio.Semaphore}-based
  concurrency gate, with a default ceiling of 20 in-flight guest AI requests;
  \item expiry after 32 hours maximum or 8 hours of inactivity.
\end{itemize}

Cleanup is part of the design, not an afterthought. The
\texttt{guest-sandbox-cleanup} Edge Function and related SQL cleanup paths release quota counters,
delete guest storage objects and sandbox-scoped class data, and remove anonymous auth users so that
demo sessions do not leak storage, quota state, or stale identities over time.

This design gives the platform a much stronger demonstration surface than a read-only walkthrough
would provide, while still keeping safety, resource control, and lifecycle hygiene under explicit
system control.

\section{Software Development Process}

\subsection{Requirements and scope shaping}

The initial project brief pushed the team toward a stakeholder-aware educational system rather than
toward a shallow AI showcase. The resulting product scope is therefore not a universal learning
platform or a generic study chatbot. It is a structured STEM classroom assistant built around
teacher materials, teacher editorial control, and student-facing learning workflows grounded in the
same class knowledge.

A practical strength of the delivery was scope discipline. The team did not try to implement every
possible education feature. Instead, it converged on one coherent classroom operating loop:
\begin{itemize}
  \item content ingestion;
  \item blueprint generation and publishing;
  \item student chat and assignments;
  \item review and analytics;
  \item guest evaluation flow.
\end{itemize}

That restraint is important because it concentrated engineering effort on one integrated product
system instead of spreading effort across unrelated AI experiments.

\subsection{Design process}

The design process was iterative, documentation-backed, and architecture-aware. The project
maintained:
\begin{itemize}
  \item top-level design, architecture, UI/UX, and deployment documents;
  \item internal reference notes under \texttt{.codex/};
  \item Mermaid diagrams exported to PNG for communication and reporting;
  \item commit history showing multiple phases of refinement rather than one-shot implementation.
\end{itemize}

UI and UX were also treated as a first-class engineering workstream rather than as late cosmetic
polish. The development record shows repeated investment in design tokens, navigation restructuring,
shared auth-surface redesign, teacher-versus-student flow separation, and consistency hardening.
That matters in an educational product because quality depends not only on AI correctness, but also
on clarity, trust, and ease of use in everyday classroom operation.

\subsection{Implementation phases observed in git history}

The git history shows staged delivery rather than a single end-loaded implementation push.
Table~\ref{tab:timeline} summarizes the major development phases visible across the repository.

\begin{table}[H]
\centering
\caption{Major development phases observed from git history}
\label{tab:timeline}
\begin{tabularx}{\textwidth}{>{\bfseries}p{2.8cm}X}
\toprule
Period & Main progression visible in the repository \\
\midrule
2026-02-03 to 2026-02-05 & Repository foundation, baseline schema and RLS, initial class flows, materials ingestion, provider adapters, blueprint generation, and first UI scaffolding. \\
2026-02-06 to 2026-02-10 & End-to-end vertical slices for grounded chat, quiz generation, flashcards, auth hardening, and always-on class chat workspace. \\
2026-02-22 to 2026-03-02 & Material-worker stabilization, queue-backed processing, live material readiness feedback, class-shell UX improvements, design-token refactoring, and stronger web test coverage. \\
2026-03-10 to 2026-03-15 & Python-backend migration and hardening, canonical backend envelope enforcement, service-auth tightening, materials file management, class intelligence delivery, and architecture/documentation alignment. \\
2026-03-18 to 2026-03-29 & Generative canvas, adaptive teaching brief, review-surface improvements, guest mode rollout, E2E expansion, shared auth redesign, branded email flows, and OTP-based password change. \\
2026-03-30 to 2026-04-05 & Inline code documentation, analytics fixes, guest quota redesign, concurrency hardening, environment/document refresh, diagram alignment, and final report refinement. \\
\bottomrule
\end{tabularx}
\end{table}

The overall delivery pattern is clear: establish the domain model, build the main classroom slices,
consolidate the backend and runtime boundaries, expand teacher-support and guest-demo capabilities,
and then harden the system, documentation, and submission artifacts near the end of the project.

\subsection{Collaboration workflow}

The collaboration model combined clear primary technical leadership with meaningful supporting
contributions across the rest of Group \#\GroupNumber. Haochen Fa served as the primary engineer
and solution architect, leading the overall technical direction, system architecture, cross-service
integration, deployment alignment, and most feature delivery. Other team members contributed through
code review, feature validation, edge-case testing, product feedback, documentation work, and
selected implementation tasks.

GitHub was the formal engineering workspace for the project. The team used GitHub Issues and pull
requests to track implementation work, review changes, and communicate technical decisions, while
WhatsApp was used for broader day-to-day discussion, rapid coordination, and design feedback. The
main development repository is \PrimaryRepoUrl, and the organization repository prepared for
review and submission is \SubmissionRepoUrl.

The engineering record reinforces this workflow. The project history now spans more than 400
commits, including 92 merge commits, nine tagged release checkpoints from \texttt{v1.0.0} through
\texttt{v1.4.0}, and repeated PR-based integrations through at least PR \#102. Those merges cover
major workstreams such as guest mode, auth redesign, UI/UX improvements, analytics, documentation
refreshes, email flows, and final runtime hardening. This is the profile of a staged collaborative
build with regular integration checkpoints, not a single large delivery at the end.

\subsection{Use of GenAI during development}

GenAI-assisted development was used deliberately as part of the delivery strategy rather than as a
substitute for engineering ownership. Codex and Claude Code were used to accelerate scaffolding,
refactoring, debugging, documentation maintenance, test-support work, and repetitive implementation
tasks across the web, backend, and Supabase surfaces. The pull-request workflow also incorporated
GitHub Copilot review as an additional automated review layer.

These tools were effective because they operated under disciplined human direction. Clear
constraints, architectural context, repo-specific instructions, and fast technical validation turned
them into productive accelerators. Without that experienced human harness, the same tools would not
have produced work at the same level of alignment, efficiency, or reliability.

Generated suggestions were therefore always paired with human review. Outputs were checked against
feature behavior, code structure, testing obligations, Git and GitHub workflow hygiene,
dependency-tooling correctness, and the runtime requirements of the Next.js web app, the FastAPI
backend, and the Supabase layer. In that sense, prompt engineering and AI-tool orchestration became
part of the team's engineering process, while responsibility for correctness, maintainability, and
product quality remained fully human.

\section{Testing, Quality Assurance, and Deployment}

\subsection{Testing strategy}

The platform uses a layered testing strategy rather than relying on a single proof mechanism.

\begin{table}[H]
\centering
\caption{Current testing layers in the repository}
\begin{tabularx}{\textwidth}{>{\bfseries}p{3.3cm}X>{\raggedright\arraybackslash}p{2.7cm}}
\toprule
Layer & What it covers & Evidence in repo \\
\midrule
Web unit and integration tests & Server actions, auth flows, route guards, material actions, blueprint logic, chat flows, guest helpers, UI primitives, and validation helpers & 60 test files under \texttt{web/src} \\
Backend unit tests & Providers, schemas, class flows, blueprint generation, quiz, flashcards, materials, analytics, chat workspace, and guest rate limiting & 17 Python test files under \texttt{backend/tests} \\
End-to-end tests & Teacher navigation, student navigation, login, join class, settings, sign-out, and guest entry against a deployed environment & 9 Playwright specs under \texttt{tests/e2e} \\
\bottomrule
\end{tabularx}
\end{table}

Representative quality signals include:
\begin{itemize}
  \item guest middleware tests plus guest sandbox and quota tests;
  \item blueprint action tests;
  \item materials extraction, preview, retrieval, and action tests;
  \item chat generation, compaction, and workspace tests;
  \item analytics, insights, and teaching brief action tests;
  \item Playwright flows across teacher, student, settings, auth sign-out, and guest entry.
\end{itemize}

\subsection{Why the testing approach is appropriate}

This test distribution matches the system's runtime boundaries:
\begin{itemize}
  \item unit tests catch domain logic regressions close to the code;
  \item web integration tests validate server actions and route behavior;
  \item backend tests verify AI-related orchestration and guardrails without needing full browser runs;
  \item Playwright tests confirm that the deployed user flows still work end to end.
\end{itemize}

For a system that spans a Next.js application, a separate FastAPI service, hosted Supabase state,
and queue-backed workers, this layered model is the right testing shape. It contains regressions at
the domain layer, verifies critical server-side workflow behavior, and still proves that the hosted
teacher, student, and guest journeys function correctly after deployment.

\subsection{Deployment topology}

The hosted platform is deployed as three coordinated surfaces: a frontend deployment, a separate
Python backend deployment, and a hosted Supabase project. Figure~\ref{fig:deployment-topology}
captures this structure.

\begin{figure}[H]
  \centering
  \FullPageWidthGraphic{DEPLOYMENT-diagram-01.png}
  \caption{Hosted deployment topology across frontend, backend, and Supabase services.}
  \label{fig:deployment-topology}
\end{figure}

This topology gives the product clearer operational boundaries:
\begin{itemize}
  \item frontend behavior can be separated from backend or data-plane faults during debugging;
  \item AI orchestration can evolve independently from the web deployment surface;
  \item Supabase-specific responsibilities such as migrations, Edge Functions, and secrets remain explicit;
  \item preview and staging environments can be aligned more carefully across the stack.
\end{itemize}

The deployment runbook also makes clear that this is not a nominal split. Both \texttt{web/} and
\texttt{backend/} carry Vercel project linkage, while Supabase hosts authentication, PostgreSQL,
storage, RLS policies, SQL dispatch helpers, and the \texttt{material-worker} and
\texttt{guest-sandbox-cleanup} Edge Functions. In other words, the deployed system is genuinely
multi-surface, not a frontend-plus-database simplification.

\subsection{Deployment process}

The documented deployment workflow includes:
\begin{enumerate}
  \item configuring one Supabase project per environment;
  \item applying the active migration set;
  \item deploying the \texttt{material-worker} and \texttt{guest-sandbox-cleanup} Edge Functions;
  \item configuring Edge Function secrets plus the Vault-backed values used by SQL-side dispatch;
  \item deploying the Python backend;
  \item deploying the web application;
  \item running post-deploy smoke tests for auth, teacher flows, student flows, insights, and guest mode.
\end{enumerate}

Operationally, this matters because deployment is not just a frontend publish. The schema, data
plane, background workers, backend environment contract, and post-deploy smoke checks all have to
line up before the product is truly usable. The material pipeline, in particular, depends on queue
dispatch, worker execution, embedding configuration, and secret placement across multiple surfaces.

Guest mode makes that need for alignment even more visible. The hosted environment has to keep the
Supabase quota RPCs, guest schema, backend guest-limit environment variables, anonymous-auth
configuration, and cleanup worker synchronized so that session admission, inactivity expiry, and
AI-concurrency guardrails behave consistently across the stack.

\subsection{Database evolution}

The active schema history now contains 23 migrations, which documents iterative platform evolution
rather than a static schema frozen after the first implementation pass.

Major migration themes include:
\begin{itemize}
  \item initial schema and RLS foundation;
  \item quiz integrity and auth hardening;
  \item canonical blueprint snapshots and publication hardening;
  \item always-on class chat and context compaction;
  \item queue-backed material processing;
  \item class insights and teaching brief snapshots;
  \item guest mode schema, seed data, cleanup, and quota redesign.
\end{itemize}

This migration history matters because it records the product's growth in a traceable and reviewable
form. The database changes do not read like generic table churn; they align closely with major
capability additions such as blueprint governance, persistent chat, async material processing,
teacher intelligence, and sandboxed guest evaluation.

\section{Lessons Learned and Reflection}

\subsection{Challenges that sharpened the product}

The most useful lesson from the development record is that the hardest parts of the platform were
not always the most visible ones. On the surface, the system looks like a clean teacher-student
learning product with AI assistance. In the commit history, however, the real engineering effort is
visible in the repeated cycles of correction, hardening, and redesign that turned that product idea
into something dependable. The guest-mode work is the clearest example. What appears in the UI as a
simple ``continue as guest'' action ultimately required anonymous authentication, sandbox cloning,
quota governance, expiry handling, cleanup flows, backend ownership checks, and later a redesign
from per-IP entry limiting to a global session-quota model. That progression is important because it
shows that the team did not stop at a demo-friendly implementation; it kept refining the subsystem
until it behaved like a controlled production feature.

Authentication and account recovery taught a similar lesson about trust. The shared auth surface now
looks cohesive, but the repository history shows that getting there required several iterations:
home-first auth redesign, resend and confirmation recovery flows, guest-signup loop fixes, branded
email templates, and OTP-based password reauthentication for sensitive settings changes. The lesson
here is that user trust is built through a chain of small reliability decisions, not through one
large feature. In an education platform, where teachers and students depend on predictable access,
these flows are part of the core product, not peripheral polish.

Another lesson was that multi-surface deployment demands operational discipline very early. Once the
platform was split across the Next.js web app, the FastAPI backend, Supabase storage and database,
and two Edge Functions, configuration mistakes stopped being minor inconveniences and became real
delivery risks. The repository's maintained notes about \texttt{trust\_env=False}, Edge Function
secret placement, queue-worker configuration, and environment alignment are not incidental trivia.
They record issues that only become visible when a system is actually being deployed and debugged
across several runtime boundaries. The team effectively learned that architecture is only as strong
as its operational contract.

Finally, the project reinforced that AI features become dependable only when they are normalized
into product workflows. Blueprint generation, analytics narratives, teaching briefs, and chat
responses could not remain raw model outputs. They had to be validated, cached, persisted,
versioned, or constrained by class context before they were trustworthy enough for real use. The
analytics fixes, teaching-brief cleanup, chat compaction logic, and blueprint publication controls
all point to the same lesson: GenAI becomes valuable in a learning system only after software
engineering disciplines it into repeatable application behavior.

\subsection{Advantages proven during development}

The strongest advantage proven during development is the product's coherence. The repository did not
drift into a loose collection of unrelated AI experiments because the Course Blueprint stayed at the
center of the design. Materials ingestion, blueprint editing, activity generation, student chat,
review, and teacher insight surfaces all continue to connect back to the same instructional object.
That coherence is visible not only in the current codebase but also in the sequence of feature
delivery: new capabilities were repeatedly added by extending the blueprint-centered learning loop
rather than by bypassing it.

A second proven advantage is the architecture boundary between interface, orchestration, and data
surfaces. During development, the team was able to redesign auth, harden guest access, improve
analytics, extend teaching-brief logic, and refine material processing without collapsing the entire
system into one tangled codepath. The dedicated Python backend, the queue-backed material worker,
and the Supabase policy and migration layer gave the product room to evolve. That separation reduced
coupling and made later refinements more manageable, especially once the project moved beyond early
vertical slices into more operational work.

The development record also shows a clear advantage in teacher-centered product thinking. Many
student-facing educational demos become generic once they add AI, but this project repeatedly
strengthened the instructor's side of the system instead: editable blueprint governance, assignment
review flows, a teacher chat monitor, adaptive teaching briefs, class intelligence snapshots, and
student-preview mode inside the same class shell. That matters because it gives the platform a
practical classroom identity. The product is not merely helping students ask questions; it is also
helping teachers prepare, supervise, interpret, and intervene.

Another advantage that became more convincing over time is demonstrability. Guest mode is not a
static showcase or a scripted walkthrough. It is a real sandbox that clones seeded class data,
supports role switching, and cleans up after itself. That is a strategic advantage for an academic
project because it makes the system easier to evaluate honestly. A reviewer can interact with the
same core flows that real users would see, while the system still preserves data isolation and
resource control.

Finally, the project benefited from disciplined documentation and verification. The current repo
contains aligned design, architecture, deployment, UI, and subsystem guides, plus layered web,
backend, and end-to-end tests. The commit history shows that documentation was updated alongside
major changes rather than being postponed until submission week. That habit reduced drift between
what the system claimed to do and what it actually did, which is one of the clearest signs of
engineering maturity in the repository.

\subsection{Reasonable next-stage directions}

The most practical next step for the platform is to connect insight more directly to follow-up
teaching action. The current system can already identify class weaknesses, summarize misconceptions,
and recommend next steps through the teaching brief and class intelligence dashboard. A sensible
extension would be to let those recommendations lead into teacher-reviewable follow-up work such as
targeted quizzes, intervention flashcard sets, or focused chat assignments tied to weaker topics.
Because the platform already has blueprint grounding, analytics snapshots, assignment flows, and
editable generation surfaces, this would extend the current design rather than introduce a new
product direction.

Another reasonable direction is to broaden analytics from a single-class view into a more
longitudinal teaching support view. At present, the system is already capable of summarizing one
class coherently. A later iteration could help teachers compare patterns across classes, track
recurring misconceptions over time, and relate assigned activities more clearly to later student
improvement. This would remain consistent with the current teacher-support focus of the platform.

The third direction is operational rather than pedagogical: stronger visibility into system health.
The platform now includes provider fallback, queue-backed workers, Edge Functions, sandbox quotas,
background cleanup, and multi-surface deployment concerns. Over time, it would be useful to expose
worker status, guest-capacity pressure, provider failures, and environment alignment checks more
explicitly. That would not change the classroom workflow directly, but it would make the system
easier to operate and maintain as the project continues to mature.

\section{Conclusion}

This project was not designed as a generic AI study tool with classroom branding added afterwards.
It was designed around a more specific educational idea: teachers should be able to turn their own
materials into a Course Blueprint, keep editorial control over what is published, and then use that
blueprint to guide what students learn, what the AI can generate, and what the teacher can later
review. That choice gave the platform a coherent center. Materials, activities, chat, review,
teaching brief, and analytics all connect back to the same class-specific instructional object.

The current repository is also stronger than a one-path prototype because the implementation
supports that idea at system level. The platform combines a dedicated Python backend for AI
orchestration, queue-backed material processing through Supabase, policy-enforced data access,
teacher-facing insight snapshots, and a guest sandbox model that allows realistic evaluation without
using real classroom data. These choices add complexity, but they also make the project more
credible as a working software system.

The development process also showed that the quality of the final result depended heavily on
hardening and revision. Guest mode needed redesign and quota control. Authentication and recovery
needed repeated refinement. Analytics, chat, and material flows needed to be normalized so they
behaved like stable product features rather than raw AI outputs. This was an important part of the
project outcome: progress came not only from adding features, but from improving reliability,
consistency, and maintainability.

Taken together, the platform represents a serious and well-scoped course group project. It has a
clear educational purpose, a credible technical architecture, meaningful GenAI integration, and a
development record that shows real attention to testing, documentation, and deployment practice.
The current implementation also leaves room for careful future improvement without overstating what
has already been achieved.

\newpage
\section{User and Operations Manual}

\noindent This section is written as a practical operating guide rather than a feature inventory. Each role
starts with a quick path to first success, then moves into compact reference material for the tasks
that are used most often during deployment, teaching, and study.

\begin{table}[H]
\centering
\footnotesize
\begin{tabularx}{\textwidth}{>{\bfseries\raggedright\arraybackslash}p{2.7cm}>{\raggedright\arraybackslash}p{2.3cm}X>{\raggedright\arraybackslash}p{3.2cm}}
\toprule
Role & Start here & Primary outcome & Main surfaces \\
\midrule
Administrator & Supabase project setup plus hosted web/backend environments & Healthy deployment, aligned configuration, and first-line operational support & Supabase Dashboard, SQL Editor, Edge Functions, web deployment, backend deployment \\
Instructor & Landing-page auth modal or \texttt{/register} & Publish a blueprint-grounded class and assign the first learning activity & Teacher Dashboard, class workspace, Blueprint Studio, review pages, Insights \\
Student & Landing-page auth modal or \texttt{/register} & Join a class, study the published blueprint, and complete assigned work & Student Dashboard, \texttt{/join}, class workspace, Settings, Help \\
\bottomrule
\end{tabularx}
\end{table}

\begin{figure}[H]
  \centering
  \includegraphics[
    width=0.58\textwidth,
    keepaspectratio
  ]{UIUX-diagram-02.png}
  \caption{Role-aware navigation from public entry to dashboards and the shared class workspace shell.}
\end{figure}

\subsection{Administrator Manual (Deployment and Operations)}

\noindent The administrator role owns deployment alignment across four surfaces: the Next.js web
application, the standalone Python backend, Supabase itself, and the two required Edge Functions.
The practical goal is simple: keep those surfaces configured with matching URLs, matching service
keys, and a working post-deploy smoke test.

\subsubsection*{Quick Start}

\begin{table}[H]
\centering
\footnotesize
\begin{tabularx}{\textwidth}{>{\bfseries\raggedright\arraybackslash}p{1.1cm}>{\raggedright\arraybackslash}p{5.9cm}X}
\toprule
Step & Do this & Done when \\
\midrule
1 & Create the target Supabase project and enable email/password auth, email confirmation, and Anonymous Auth if guest mode is required. & Auth settings and redirect URLs match the intended environment. \\
2 & Apply database migrations from \texttt{supabase/migrations}. & Core tables, RLS, guest-mode schema, and support RPCs exist. \\
3 & Deploy \texttt{material-worker} and \texttt{guest-sandbox-cleanup}. & Both functions appear in Supabase and have their required secrets. \\
4 & Configure web and backend environment variables, especially \texttt{PYTHON\_BACKEND\_URL} and \texttt{PYTHON\_BACKEND\_API\_KEY}. & Web requests can authenticate to the backend successfully. \\
5 & Deploy the web app and the Python backend as separate hosted services. & The live site loads and \texttt{/healthz} responds from the backend. \\
6 & Run a smoke test: sign in, open a dashboard, upload one material, and confirm AI-backed surfaces can respond. & The release is operational end to end, not just visually deployed. \\
\bottomrule
\end{tabularx}
\end{table}

\subsubsection*{Essential configuration}

\begin{table}[H]
\centering
\footnotesize
\begin{tabularx}{\textwidth}{>{\bfseries\raggedright\arraybackslash}p{2.3cm}>{\raggedright\arraybackslash}p{5.5cm}X}
\toprule
Surface & Essential variables & Why they matter \\
\midrule
Web app & Supabase URL/key, site URL, admin secret, backend URL/key & Enables auth redirects, admin server actions, and authenticated calls to the backend. \\
Backend & Shared backend key, one AI provider, Supabase service key, worker token & Allows the backend to serve AI routes, talk to Supabase, and dispatch material jobs. \\
Material worker & Supabase service key, embedding provider secret, worker token & Required for queue-backed chunking, embedding, and material status updates. \\
Guest cleanup function & Supabase service key, cleanup token & Required to expire guest sandboxes, release quota, and remove demo artifacts safely. \\
\bottomrule
\end{tabularx}
\end{table}

\textbf{Minimal web example (\texttt{web/.env.local})}

\begin{verbatim}
NEXT_PUBLIC_SUPABASE_URL=https://<project-ref>.supabase.co
NEXT_PUBLIC_SUPABASE_PUBLISHABLE_KEY=<publishable-key>
NEXT_PUBLIC_SITE_URL=http://localhost:3000
SUPABASE_SECRET_KEY=<service-role-key>
PYTHON_BACKEND_URL=http://localhost:8001
PYTHON_BACKEND_API_KEY=<shared-backend-key>
NEXT_PUBLIC_GUEST_MODE_ENABLED=false
\end{verbatim}

\textbf{Minimal backend shell or env-loader example}

\begin{verbatim}
PYTHON_BACKEND_API_KEY=<same-shared-backend-key>
AI_PROVIDER_DEFAULT=openrouter
OPENROUTER_API_KEY=<provider-key>
OPENROUTER_MODEL=<chat-model>
OPENROUTER_EMBEDDING_MODEL=<embedding-model>
SUPABASE_URL=https://<project-ref>.supabase.co
SUPABASE_SERVICE_ROLE_KEY=<service-role-key>
MATERIAL_WORKER_TOKEN=<edge-function-token>
\end{verbatim}

\textbf{Edge Function notes}
\begin{itemize}
  \item Keep Edge Function secrets in Supabase Edge Function Secrets, not in Vercel-only configuration.
  \item Apply the branded confirmation and recovery templates in Supabase Auth before stakeholder demos.
  \item When guest mode is enabled, treat Anonymous Auth, guest seed data, and cleanup automation as one deployment feature.
\end{itemize}

\subsubsection*{Core commands}

\begin{verbatim}
# Local setup and verification
pnpm install
pip install -r backend/requirements.txt
pnpm dev
uvicorn app.main:app --app-dir backend --host 0.0.0.0 --port 8001 --reload
pnpm test
python3 -m unittest discover -s backend/tests -p 'test_*.py'

# Supabase deploy and function updates
npx supabase db push
npx supabase functions deploy material-worker
npx supabase functions deploy guest-sandbox-cleanup

# Frontend deploy and backend health
cd web
npx --yes vercel --yes --prod
curl https://<backend-domain>/healthz
\end{verbatim}

\subsubsection*{Frequent support reference}

\begingroup
\small
\setlength\LTleft{0pt}
\setlength\LTright{0pt}
\begin{longtable}{@{}>{\bfseries\raggedright\arraybackslash}p{2.6cm}>{\raggedright\arraybackslash}p{5.1cm}>{\raggedright\arraybackslash}p{6.0cm}@{}}
\toprule
Task & Query or action & Use this when \\
\midrule
\endfirsthead
\toprule
Task & Query or action & Use this when \\
\midrule
\endhead
\bottomrule
\endfoot
\bottomrule
\endlastfoot
Review accounts & Join \texttt{auth.users} with \texttt{public.profiles}. & You need to confirm email state, display name, or whether a user is teacher or student. \\
Register admin & Insert the user's UUID into \texttt{public.admin\_users}. & A new operator needs platform-admin capability without changing the normal profile role model. \\
Review classes & Query \texttt{public.classes} for title, subject, join code, archived state. & You need to verify class creation, join-code sharing, or archive status. \\
Inspect materials & Join the materials table with the processing-job table. & A teacher reports a material stuck in \texttt{processing} or \texttt{failed}. \\
Inspect guest sandboxes & Query \texttt{public.guest\_sandboxes} for role, status, expiry, and last activity. & Guest-mode cleanup or quota behaviour looks wrong. \\
Manual guest cleanup & Run the guest cleanup RPC from the SQL Editor. & Cleanup is lagging and you need a manual pass from the SQL Editor. \\
\end{longtable}
\endgroup

\noindent In practice, the two SQL views worth saving in the Supabase editor are the
\texttt{auth.users + public.profiles} join for account support and the
\texttt{materials + material\_processing\_jobs} join for stuck uploads.

\subsubsection*{Troubleshooting}

\begin{table}[H]
\centering
\small
\begin{tabularx}{\textwidth}{>{\bfseries\raggedright\arraybackslash}p{3.2cm}>{\raggedright\arraybackslash}p{4.7cm}X}
\toprule
Symptom & First checks & Likely recovery action \\
\midrule
Guest entry fails & Guest-mode web flag, Anonymous Auth, guest seed data, cleanup function deployment & Re-enable Anonymous Auth, confirm guest schema/seed state, and verify cleanup-tokenized deployment. \\
Material stays in \texttt{processing} & \texttt{material-worker} deployment, worker token, job \texttt{attempts} and \texttt{last\_error}, embedding provider secrets & Redeploy the worker, fix missing secrets, then re-run the teacher upload path. \\
Auth links use the wrong domain & \texttt{NEXT\_PUBLIC\_SITE\_URL}, Supabase Site URL, redirect URL list & Align local, preview, and production URLs so confirmation and recovery callbacks return to the right environment. \\
Web UI loads but AI actions fail & Backend \texttt{/healthz}, \texttt{PYTHON\_BACKEND\_URL}, shared backend API key, provider keys & Restore backend reachability first, then verify provider and Supabase service credentials. \\
\bottomrule
\end{tabularx}
\end{table}

\subsection{Instructor Manual}

\noindent The instructor workflow is anchored by the blueprint lifecycle: draft, review, approve,
publish. Students only receive grounded AI chat and blueprint-powered activities after the class
blueprint has been published, so the fastest way to productive use is to move one class cleanly
through that pipeline.

\subsubsection*{Quick Start}

\begin{enumerate}
  \item Register as a \textbf{teacher} from the landing page or \texttt{/register}, confirm the email,
  and open the \textbf{Teacher Dashboard}.
  \item Choose \textbf{Create class}, enter the class title and any optional subject or level details,
  then save the class.
  \item Open the class workspace and upload the first PDF, DOCX, or PPTX materials.
  \item Wait until uploaded materials move from \texttt{processing} to \texttt{ready}.
  \item Open \textbf{Blueprint Studio}, generate the draft, review it, approve it, and publish it.
  \item Create the first chat assignment, quiz, or flashcards activity and assign or publish it for
  students.
\end{enumerate}

\textbf{Before students can use AI-backed learning surfaces}
\begin{itemize}
  \item Materials should be uploaded and marked \texttt{ready}.
  \item The blueprint should be published, not only saved as a draft.
  \item Student-facing activities should be published or assigned from the class workspace.
\end{itemize}

\subsubsection*{Core tasks}

\begin{table}[H]
\centering
\small
\begin{tabularx}{\textwidth}{>{\bfseries\raggedright\arraybackslash}p{2.8cm}>{\raggedright\arraybackslash}p{3.9cm}X}
\toprule
Task & Where in the UI & Practical result \\
\midrule
Create class & Teacher Dashboard or \textbf{My Classes} & A new class shell is created and a join code becomes available for student enrollment. \\
Upload materials & Class workspace material area & Source files are stored, processed, and made available for preview, download, and blueprint generation. \\
Publish blueprint & \textbf{Blueprint} page in the class workspace & The class receives a published study guide that unlocks grounded AI chat and activity generation. \\
Create chat assignment & Class overview \texttt{New chat assignment} path & Students receive a guided conversation flow with persisted transcript and reflection review. \\
Generate quiz & Quiz creation route from the class page & A blueprint-grounded quiz draft can be edited, published, and assigned. \\
Generate flashcards & Flashcards creation route from the class page & A flashcard set can be refined, published, and assigned for review practice. \\
Review submissions & Assignment review pages & Completion state, answers, transcripts, scores, and teacher feedback are visible in one place. \\
Use Insights & \textbf{Insights} page for the class & Charts, summaries, AI narrative, and intervention suggestions support follow-up teaching action. \\
Preview student view & Class workspace student-preview mode & The instructor can check the learner experience without switching accounts. \\
\bottomrule
\end{tabularx}
\end{table}

\noindent Recommended sequence for a new class: set up the class, upload and verify materials,
publish the blueprint, create the first activities, then use review pages and Insights as student
work starts to arrive.

\subsubsection*{Reference notes}

\begin{itemize}
  \item Public auth uses the same shared surface as the student side. Teachers can start from the
  home modal or the fallback routes \texttt{/login}, \texttt{/register}, and
  \texttt{/forgot-password}.
  \item Student preview is useful for validating that published blueprint widgets, AI chat, and
  assignment surfaces are visible exactly as intended.
  \item The current release centers teaching activity around three implemented activity families:
  chat assignments, quizzes, and flashcards.
\end{itemize}

\subsection{Student Manual}

\noindent The student experience is intentionally narrower than the teacher experience. Its job is to
help learners join the right class quickly, understand what is available now, and move between the
published blueprint, AI chat, and assigned activities without crossing into teacher-only controls.

\subsubsection*{Quick Start}

\begin{enumerate}
  \item Register as a \textbf{student} from the landing page or \texttt{/register}, confirm the email,
  and open the \textbf{Student Dashboard}.
  \item Open \textbf{Join class} from the dashboard or sidebar and enter the join code from the
  instructor.
  \item Open the joined class workspace and check whether the blueprint and activities are already
  published.
  \item Use the published blueprint first to review topics and learning objectives.
  \item Open AI chat or the assigned quiz, flashcards, or chat assignment and complete the required
  work.
  \item Revisit the dashboard, Settings, or Help pages whenever you need to track progress or manage
  account details.
\end{enumerate}

\subsubsection*{Dashboard and assignment states}

\begin{table}[H]
\centering
\small
\begin{tabularx}{\textwidth}{>{\bfseries\raggedright\arraybackslash}p{2.8cm}>{\raggedright\arraybackslash}p{3.9cm}X}
\toprule
State or area & Where to look & What it means \\
\midrule
Due now & Student Dashboard & Work that needs immediate attention and should be opened first. \\
Upcoming & Student Dashboard & Assigned work that is visible but not yet urgent. \\
Completed & Student Dashboard & Work that has already been finished and can be revisited for review. \\
Submitted or reviewed & Inside the assignment page & The latest attempt or reflection has been saved, and teacher review may now be visible. \\
\bottomrule
\end{tabularx}
\end{table}

\subsubsection*{Tools inside the class workspace}

\begin{table}[H]
\centering
\small
\begin{tabularx}{\textwidth}{>{\bfseries\raggedright\arraybackslash}p{2.7cm}X>{\raggedright\arraybackslash}p{3.3cm}}
\toprule
Surface & What it is for & Important note \\
\midrule
Published blueprint & Study guide for topics, objectives, and approved class focus & It is the best starting point before quizzes or AI chat. \\
AI chat & Ask grounded questions about the class content & Chat is available after the instructor publishes the blueprint. \\
Chat assignment & Complete a guided conversation and submit the required reflection & This is different from open practice chat because it is reviewed as class work. \\
Quiz & Answer blueprint-grounded questions and submit the attempt & Reopen the assignment later to review score and feedback state. \\
Flashcards & Review generated cards and submit progress & Good for repetition and pre-quiz revision. \\
Settings & Update display name and use the password-change flow & Password change uses current-password verification plus a code step. \\
Help & Open role-specific FAQ and next-step guidance & Useful when you are not sure where a class feature is located. \\
\bottomrule
\end{tabularx}
\end{table}

\subsubsection*{If something seems missing}

\begin{table}[H]
\centering
\small
\begin{tabularx}{\textwidth}{>{\bfseries\raggedright\arraybackslash}p{3.1cm}>{\raggedright\arraybackslash}p{4.8cm}X}
\toprule
Situation & First thing to check & Most likely explanation \\
\midrule
No class visible & Confirm you joined successfully with the correct code from \texttt{/join}. & The class has not been joined yet, or the code was entered incorrectly. \\
AI chat is unavailable & Check whether the class blueprint is published. & AI chat is intentionally gated behind the published blueprint. \\
No assignments appear & Re-check the dashboard and the class workspace. & The instructor may not have published or assigned the activity yet. \\
Need to change account details & Open \textbf{Settings} or \textbf{Help}. & Profile rename and password management live outside the class workspace. \\
\bottomrule
\end{tabularx}
\end{table}

\subsection{Supplementary Guest Demo Flow}

Guest mode is not part of the main three-role operating model, but it is important for demos and
evaluation. It behaves like a sandboxed temporary classroom rather than a permanent account.

\begin{table}[H]
\centering
\small
\begin{tabularx}{\textwidth}{>{\bfseries\raggedright\arraybackslash}p{2.7cm}>{\raggedright\arraybackslash}p{3.8cm}X}
\toprule
Action & Where it starts & What happens \\
\midrule
Enter guest mode & Landing page \textbf{Continue as guest} button when guest mode is enabled & The app creates an anonymous authenticated session and provisions a sandbox classroom. \\
Explore both roles & Guest class workspace & The demo can switch between teacher and student perspectives without affecting real user data. \\
Reset sandbox & Guest controls inside the sandbox flow & A fresh walkthrough can be created when the current one becomes cluttered. \\
Create permanent account & Exit the guest-only path and use normal auth & Demo exploration stays isolated from real teacher or student accounts. \\
\bottomrule
\end{tabularx}
\end{table}

\newpage
\appendix

\section{Repository Evidence Used for This Report}

This report was prepared against the current repository state using the following main sources:

\begin{itemize}
  \item top-level documents:
  \begin{itemize}
    \item \texttt{README.md}
    \item \texttt{ARCHITECTURE.md}
    \item \texttt{DESIGN.md}
    \item \texttt{UIUX.md}
    \item \texttt{DEPLOYMENT.md}
    \item \texttt{AGENTS.md}
    \item \texttt{CLAUDE.md}
  \end{itemize}
  \item subsystem documentation:
  \begin{itemize}
    \item \texttt{web/README.md}
    \item \texttt{backend/README.md}
    \item \texttt{supabase/README.md}
  \end{itemize}
  \item course materials:
  \begin{itemize}
    \item \texttt{project\_overivew.pdf}
    \item \texttt{project\_background.pdf}
  \end{itemize}
  \item implementation evidence:
  \begin{itemize}
    \item key web routes and server actions
    \item backend service modules
    \item Supabase Edge Functions
    \item migration history
    \item git commit history and merge history
    \item test suites and E2E specs
  \end{itemize}
\end{itemize}

\section{Submission and Live System Links}

\begin{itemize}
  \item Live product: \LiveProductUrl
  \item Demo video: \DemoVideoUrl
  \item Primary development repository: \PrimaryRepoUrl
  \item Organization submission repository: \SubmissionRepoUrl
\end{itemize}

\end{document}
